% Shared preamble for the two point-by-point response letters (one per reviewer).
\documentclass[11pt,a4paper]{article}
\usepackage[margin=2.4cm]{geometry}
\usepackage{amsmath,amssymb}
\usepackage[T1]{fontenc}
\usepackage{lmodern}
\usepackage{xcolor}
\usepackage{framed}
\usepackage{enumitem}
\usepackage[hidelinks]{hyperref}
\DeclareMathOperator{\spec}{spec}
\DeclareMathOperator{\diag}{diag}
\setlength{\parskip}{4pt}
\setlength{\parindent}{0pt}
\definecolor{shadecolor}{gray}{0.93}
% \comment{label}{text}: the reviewer's comment, quoted in a shaded box.
\newcommand{\comment}[2]{\par\medskip\begin{shaded}\textbf{Comment #1.} \emph{#2}\end{shaded}}
% \response{text}
\newcommand{\response}[1]{\par\textbf{Response.} #1\par\medskip}
% \where{text}: pointer into the revised manuscript (section, equation, page/line numbers of the revised PDF).
\newcommand{\where}[1]{\par\textit{Location in the revised manuscript:} #1\par}
\newcommand{\manuscriptid}{mathematics-4633263}
\newcommand{\manuscripttitle}{From Cain to Matignon: Exact Positive-Diagonal Stability in Dimension Three and a Nonlinear Caputo Ecological Realization}
\newcommand{\letterhead}[1]{%
  \begin{center}
  {\Large\bfseries Response to #1}\\[6pt]
  Manuscript \texttt{\manuscriptid}, \emph{Mathematics} (MDPI)\\[2pt]
  \emph{\manuscripttitle}
  \end{center}
  \medskip
  Page and line numbers below refer to the revised manuscript (line numbers printed in the margin of the revised PDF). Every change requested by #1 is addressed individually; changes made in response to the other reviewer are listed at the end for completeness, because they affect the same sections.\par\medskip}

\begin{document}
\letterhead{Reviewer 1}

We thank the reviewer for the careful reading and for the precise suggestions, all of which were adopted. The reviewer's summary was correct: none of the seven points required new analysis or new numerical experiments, and the revision for this reviewer is confined to making implicit steps explicit, aligning precision statements with what is printed, and removing internal labels.

\comment{1}{Make the cubic boundary factorization explicit. In Lemma 1, display the factorization \(p(\lambda)=(\lambda^2-2ur\lambda+r^2)(\lambda+a+2ur)\), \(u=\cos\theta\), where \((1-4u^2)r^2-2uar-b=0\) and \(c=r^2(a+2ur)\), which identifies the boundary roots \(\lambda_{1,2}=re^{\pm i\theta}\) and \(\lambda_3=-a-2ur<0\).}
\response{Added exactly as suggested. The proof of Lemma 1 now displays the factorization, states the quadratic relation \((1-4u^2)r^2-2uar-b=0\) and the identity \(c=r^2(a+2ur)\), and identifies the boundary conjugate pair \(re^{\pm i\theta}\) and the strictly negative third root \(-a-2ur\).}
\where{Section 4.1, proof of Lemma 1, page 7, lines 226--232.}

\comment{2}{Clarify attainment of the threshold minimum and its classical endpoint: note that \(h_\alpha(0)>0\) for \(2/3<\alpha<1\), so that \(h_\alpha(B_\beta(x))/(x_1x_2x_3)\to+\infty\) at \(\partial\Delta^2\); at \(\alpha=1\) use \(B_\beta(x)/(x_1x_2x_3)=\beta_{23}/x_1+\beta_{13}/x_2+\beta_{12}/x_3\to+\infty\) since \(h_1(0)=0\); state that \(T_1\) is the endpoint obtained from \(h_1(b)=b\).}
\response{Added. The text following the definition of \(T_\alpha\) now contains the explicit value \(h_\alpha(0)=\bigl(\tfrac{2u}{1-4u^2}\bigr)^2\bigl(1+\tfrac{4u^2}{1-4u^2}\bigr)>0\) for \(2/3<\alpha<1\), the resulting divergence of the objective at the boundary of the simplex, and the separate argument for \(\alpha=1\) through the reciprocal-coordinate expression. It states explicitly that \(T_1\) is the classical endpoint with \(h_1(b)=b\).}
\where{Section 4.2, page 8, lines 245--250, and the statement \(h_1(b)=b\) at page 9, line 301.}

\comment{3}{Add one line explaining strictness after minimization in Theorem 5: for \(2/3<\alpha_1<\alpha_2\le1\), with \(x^*\) the minimizer for \(\alpha_1\), \(T_{\alpha_1}(\beta)=h_{\alpha_1}(B_\beta(x^*))/(x_1^*x_2^*x_3^*)>h_{\alpha_2}(B_\beta(x^*))/(x_1^*x_2^*x_3^*)\ge T_{\alpha_2}(\beta)\); the same argument gives \(T_\alpha(\beta)>T_1(\beta)\).}
\response{Added. The proof of Theorem 5 now evaluates both objectives at the minimizer attained for the smaller order and displays the three-term inequality; the comparison with \(T_1\) is made by the same argument.}
\where{Section 5.2, proof of Theorem 5, page 13, lines 412--418 (displayed inequality at line 416).}

\comment{4}{Explain explicitly why the low-order endpoint \(\alpha=2/3\) is included: with \(\mu\in\spec(-DA)\) and \(\lambda=-\mu\), Kellogg's strict wedge gives \(|\arg\mu|<2\pi/3\), hence \(|\arg\lambda|=\pi-|\arg\mu|>\pi/3\ge\alpha\pi/2\) for \(0<\alpha\le2/3\).}
\response{Added in the proof of Corollary 1, in the two-line form proposed by the reviewer. The text now says explicitly that the endpoint is included because the eigenvalue inequality is strict, not by a limiting argument.}
\where{Section 4.3, proof of Corollary 1, page 10, lines 325--330 (displayed chain at line 329).}

\comment{5}{Distinguish the fractional-only set from its interior: on the strict-\(P(-A)\) stratum and for \(2/3<\alpha<1\), \(A\in\mathcal F_\alpha^{(3)}\setminus\mathcal D_H^{(3)}\iff T_1(\beta)\le\kappa<T_\alpha(\beta)\), whereas the full-dimensional interior is \(T_1(\beta)<\kappa<T_\alpha(\beta)\); at \(\kappa=T_1(\beta)\) a positive diagonal scaling produces an imaginary conjugate pair.}
\response{Added as a remark. The revised text states the exact set with the non-strict lower inequality, states that the interior requires both inequalities strict, and explains that at \(\kappa=T_1\) one positive diagonal scaling has a purely imaginary pair, which excludes Hurwitz stability but satisfies the Matignon condition for every \(\alpha<1\).}
\where{Section 4.3, page 10, lines 319--324 (displayed equivalence at line 321).}

\comment{6}{Emphasize the role of response-rate scaling in the nonlinear example: distinguish \(J=\diag(u^*)B\) from \(R_*J=D_*B\), \(R_*=D_*\diag(u^*)^{-1}\); failure of classical \(D\)-stability means \(\exists D\succ0:\spec(DB)\not\subset\{\operatorname{Re}z<0\}\) and does not imply that the unscaled Jacobian is unstable; say that the comparison uses a selected positive response-rate configuration.}
\response{Added. Section 7.2 now defines \(B=DF(u^*)\) as the reduced per-capita derivative and \(J=\diag(u^*)B\) as the Jacobian of the unscaled system, states that the angular search selects a member \(D_*B\) of the orbit rather than \(J\), and gives the reviewer's definition of failure of Hurwitz \(D\)-stability together with the explicit statement that it does not assert instability of \(J\). Section 7.3 states that the time-domain comparison uses the selected positive response-rate configuration \(R_*\), whose Jacobian is \(R_*J=D_*B\), and that all-diagonal Matignon stability is a separate analytic statement. In addition, in response to Reviewer 2, the unscaled configuration \(R=I\) was also integrated across all tested orders and decays at every order, which illustrates the reviewer's point numerically.}
\where{Section 7.2, page 24, lines 688--692; Section 7.3, page 25, lines 705--710 and the paragraph at lines 742 ff.}

\comment{7}{Align numerical precision and certification terminology with the information provided: write \(m_0\approx0.374451281242515\) in the main text and specify the precision of the high-precision computation; in Section 7.1 describe the interval software, working precision and the procedure enclosing the minimizer, in particular how a validated lower bound \(\underline T_\alpha\) with \(\underline T_\alpha-\kappa>0\) was obtained, since uniqueness of the minimizer alone does not provide it; define or remove the labels ``C-07'', ``C-09'' and ``C-10'' appearing in the text and figures.}
\response{All three parts were done.
\begin{itemize}[leftmargin=1.4em]
\item \emph{Precision.} Equation (88) now prints \(m_0\approx0.374451281242515\). The text states that the unrounded computation of the Cain-gap equation used 60 decimal digits and that the residual below \(10^{-59}\) refers to that internal approximation, not to the printed 15-digit value.
\item \emph{Certification procedure.} Section 7.1 now names the interval software (the outward-rounded \texttt{mpmath.iv} backend), the working precisions (50 decimal digits for parameter boxes and equilibrium enclosures, 40 for the threshold), and the procedure: a high-precision Newton candidate \(\widehat y\) for the minimizer of the log-objective, a subdivided box around \(\widehat y\) on which an interval lower bound \(\underline\lambda>0\) of the smallest Hessian eigenvalue is verified, the bound \(\rho_y=\overline{\|\nabla G(\widehat y)\|}/\underline\lambda\) on the displacement of the true minimizer, and the resulting validated lower bound \(\underline T_\alpha=\exp\bigl(\underline{G(\widehat y)}-\overline{\|\nabla G(\widehat y)\|}\rho_y\bigr)\le T_\alpha(\beta)\). The text states, as the reviewer requested, that uniqueness of the minimizer alone would not validate the bound, and that the certificate \(\underline T_\alpha(\underline\beta)-\overline\kappa>0\) uses the lower corner of the \(\beta\) enclosure through monotonicity of \(T_\alpha\) in each \(\beta_{ij}\).
\item \emph{Labels.} The internal labels C-07, C-09 and C-10 were removed from the text and from all captions, and the two figures whose artwork contained them (Figures 1 and 2 of the revised manuscript) were regenerated without them. The compiled revised PDF was searched and contains no occurrence of these labels.
\end{itemize}}
\where{Section 7.1, page 23, lines 639--666, Equation (88) at line 663; captions of Figures 1 and 2.}

\section*{Changes made in response to Reviewer 2 that affect the same sections}
For the reviewer's information: Section 4.3 gained a paragraph on the strata with vanishing principal minors (page 9, lines 286--297); Section 7.1 gained an interval-certified 14-parameter neighborhood of the realization (page 23, lines 667--678); Section 7.3 gained a paragraph on initial-condition perturbations, citing the linearized asymptotic stability theorem for Caputo systems (page 25, lines 718--736), and a long-horizon nonlinear integration across the exact boundary order with a new Table 2 and a new Figure 7 (pages 25--27); all figure captions were made self-contained; recent references were added to the Introduction. None of these changes alters a theorem or a proof.
\end{document}
